\documentclass[12pt]{article}

\usepackage[a4paper, margin=1in]{geometry}
\usepackage{setspace}
\usepackage{hyperref}
\usepackage{listings}
\usepackage{xcolor}
\usepackage{graphicx}

\onehalfspacing

\title{\textbf{Snake Stack Project Report}\\
A Three-Layer Web Architecture using JavaScript, Flask, and Bash}
\author{Shiny-Bhavaprita}
\date{}

\begin{document}

\maketitle

\newpage
\tableofcontents
\newpage

%--------------------------------------------------
\section{Introduction}

The Snake Stack project is a browser-based implementation of the classic Snake game designed using a three-layer architecture. The goal of this project is not only to implement a functional game but also to demonstrate clear separation of concerns across frontend, backend, and system-level scripting.

The system is composed of:
\begin{itemize}
    \item A JavaScript-based frontend for gameplay
    \item A Flask-based backend for handling requests and storing scores
    \item A Bash script for data management and analysis
\end{itemize}

Unlike typical applications, this project uses a flat file instead of a database, encouraging the use of Unix tools for processing data.

%--------------------------------------------------
\section{External Libraries Used}

\subsection{Flask}
Flask is a lightweight Python web framework used to create the backend server. It handles HTTP requests and responses.

\subsection{JavaScript Fetch API}
The Fetch API is used in the frontend to send asynchronous POST requests to the Flask server without reloading the page.

\subsection{HTML5 Canvas}
Canvas is used for rendering the Snake game. It allows drawing shapes and animations dynamically.

\subsection{Bootstrap (Optional)}
Bootstrap is used to design modals and improve UI consistency.

\subsection{Bash Utilities}
Tools such as \texttt{awk}, \texttt{sed}, and \texttt{sort} are used for parsing and analyzing stored data.

%--------------------------------------------------
\section{Features Implemented}

\subsection{Game Mechanics}
\begin{itemize}
    \item Grid-based snake movement
    \item Continuous motion system
    \item Dynamic food spawning
\end{itemize}

\subsection{Controls}
\begin{itemize}
    \item Arrow keys
    \item WASD keys
\end{itemize}

\subsection{Food Types}
\begin{itemize}
    \item Carrot (+1 score)
    \item Pumpkin Pie (+3 score)
    \item Golden Apple (temporary immunity)
\end{itemize}

\subsection{Game Over Conditions}
\begin{itemize}
    \item Collision with wall
    \item Collision with itself
\end{itemize}

\subsection{User Interface}
\begin{itemize}
    \item Start modal with username input
    \item Game over modal with statistics
    \item Clean centered layout
\end{itemize}

\subsection{Backend Features}
\begin{itemize}
    \item REST API endpoint for score submission
    \item JSON validation
    \item File-based storage
\end{itemize}

\subsection{Bash Features}
\begin{itemize}
    \item Menu-driven interface
    \item Filtering and sorting
    \item Statistical analysis
\end{itemize}

%--------------------------------------------------
\section{Three-Layer Communication}

\subsection{Frontend to Backend}
The frontend sends game results using a POST request:

\begin{lstlisting}[language=JavaScript]
fetch('/save_score', {
  method: 'POST',
  headers: {'Content-Type': 'application/json'},
  body: JSON.stringify(data)
});
\end{lstlisting}

\subsection{Backend Processing}
The backend extracts and validates data:

\begin{lstlisting}[language=Python]
@app.route('/save_score', methods=['POST'])
def save_score():
    data = request.json
    with open("history.txt", "a") as f:
        f.write(format_data(data))
    return {"status": "success"}
\end{lstlisting}

\subsection{Data Flow}
\begin{itemize}
    \item Data is written to history.txt
    \item Bash script reads file directly
\end{itemize}

This separation ensures modularity and maintainability.

%--------------------------------------------------
\section{Format of history.txt}

Example entry:

\begin{verbatim}
[2025-06-12 14:35:02] Herobrine | 42 | WALL | 85
\end{verbatim}

\subsection{Design Decisions}
\begin{itemize}
    \item Timestamp ensures chronological tracking
    \item Pipe separator simplifies parsing
    \item Human-readable format aids debugging
\end{itemize}

\subsection{Advantages}
\begin{itemize}
    \item Easy to process with Bash tools
    \item No database dependency
    \item Lightweight and portable
\end{itemize}

%--------------------------------------------------
\section{Hurdles Faced}

\subsection{Frontend Issues}
Handling smooth movement and collision detection required careful logic.

\subsection{Backend Issues}
Ensuring proper JSON validation and avoiding malformed data.

\subsection{Bash Challenges}
Parsing structured data with shell tools required learning \texttt{awk} and \texttt{sed}.

\subsection{Debugging}
Tracking bugs across three layers was complex and required systematic testing.

%--------------------------------------------------
\section{Future Improvements}

Given 10 more hours, the following enhancements could be implemented:

\begin{itemize}
    \item Leaderboard display in UI
    \item Better animations
    \item Sound effects
    \item Database integration
    \item Multiplayer mode
    \item Advanced analytics in Bash
\end{itemize}

%--------------------------------------------------
\section{Conclusion}

The Snake Stack project demonstrates a clean separation of concerns using three distinct technologies. It highlights the integration of web technologies with system-level scripting and showcases the power of modular design.

%--------------------------------------------------
\section{Bibliography}

\begin{itemize}
    \item Flask Documentation: \url{https://flask.palletsprojects.com/}
    \item MDN Web Docs: \url{https://developer.mozilla.org/}
    \item Canvas API: \url{https://developer.mozilla.org/en-US/docs/Web/API/Canvas_API}
    \item Bash Guide: \url{https://tldp.org/LDP/abs/html/}
\end{itemize}

\end{document}
